% ============================================================
%  وعده یا خدعه؟
%  Segment 4: فصل‌های ۱۳ تا ۱۶ — ارزیابی نهایی و نتیجه‌گیری
% ------------------------------------------------------------
%  ⚠ این بخش باید پس از پایان Segment 3 (قبل از \end{document})
%    الحاق شود.
%  ⚠ پریامبل تکرار نمی‌شود — فقط محتوای فصل‌ها.
% ============================================================

% ############################################################
%  بخش سوم: ارزیابی نهایی و نتیجه‌گیری
% ############################################################

\part{ارزیابی نهایی و نتیجه‌گیری}

% ============================================================
%  فصل ۱۳: تحلیل تطبیقی — خمینی در هر بزنگاه چه کرد؟
% ============================================================
\chapter{تحلیل تطبیقی: خمینی در هر بزنگاه چه کرد؟}

\begin{infoBox}[چکیدهٔ فصل]
این فصل بزنگاه‌های هفت‌گانه‌ای را که در فصل‌های ۶ تا ۱۲ بررسی شدند، در یک \keyword{ماتریس تطبیقی واحد} گرد هم می‌آورد. هدف آن است که الگوی رفتاری خمینی را در مجموع ببینیم — نه فقط در هر رویداد به‌تنهایی. برای هر بزنگاه، اقدام خمینی، پیش‌بینی هر تز، و میزان سازگاری واقعیت با آن پیش‌بینی ارزیابی می‌شود. در پایان، الگوی کلی شناسایی و تحلیل می‌شود.
\end{infoBox}

% ────────────────────────────────────────
\section{روش‌شناسی: چگونه تزها را می‌سنجیم؟}
% ────────────────────────────────────────

\begin{principleBox}[معیارهای سنجش]
برای هر بزنگاه، سه پرسش مطرح می‌شود:
\begin{enumerate}[label=\textcolor{cPrimary}{\bfseries\arabic*.}, rightmargin=2em]
    \item \textbf{اقدام خمینی چه بود؟} (واقعیت مستند)
    \item \textbf{هر تز چه پیش‌بینی می‌کرد؟} (اگر تز درست باشد، خمینی باید چه می‌کرد؟)
    \item \textbf{واقعیت با کدام پیش‌بینی بیشتر سازگار است؟} (امتیاز تطبیق: ● = بالا، ◐ = متوسط، ○ = پایین)
\end{enumerate}
\end{principleBox}

% ────────────────────────────────────────
\section{ماتریس تطبیقی اصلی}
% ────────────────────────────────────────

\begin{landscape}

\begin{longtable}{%
    >{\raggedleft\arraybackslash}p{0.11\textwidth}
    >{\raggedleft\arraybackslash}p{0.16\textwidth}
    >{\raggedleft\arraybackslash}p{0.15\textwidth}
    >{\raggedleft\arraybackslash}p{0.15\textwidth}
    >{\raggedleft\arraybackslash}p{0.15\textwidth}
    >{\raggedleft\arraybackslash}p{0.15\textwidth}}
\toprule
\rowcolor{cPrimary!15}
\textbf{بزنگاه} &
\textbf{اقدام خمینی} &
\cellcolor{cGreen!15}\textbf{تز ۱ (صداقت)} &
\cellcolor{cAccent!10}\textbf{تز ۲ (خدعه)} &
\cellcolor{cGold!15}\textbf{تز ۳ (تکامل)} &
\cellcolor{cConsolid!10}\textbf{تز ۴ (نخبگان)} \\
\midrule
\endhead

\textbf{۱. پاریس}{\footnotesize\\ (مهر–بهمن ۵۷)} &
وعدهٔ عدم حکومت، آزادی، بازگشت به قم &
\cellcolor{cGreen!8} پیش‌بینی: صادقانه \newline تطبیق: ● &
\cellcolor{cAccent!5} پیش‌بینی: تاکتیکی \newline تطبیق: ● &
\cellcolor{cGold!8} پیش‌بینی: مرحلهٔ گذار \newline تطبیق: ● &
\cellcolor{cConsolid!5} پیش‌بینی: مشاوران تلطیف \newline تطبیق: ◐ \\
\midrule

\textbf{۲. بازگشت}{\footnotesize\\ (بهمن ۵۷)} &
«من دولت تعیین می‌کنم» + انتصاب بازرگان &
\cellcolor{cGreen!8} پیش‌بینی: انتصاب غیرروحانی \newline تطبیق: ◐ &
\cellcolor{cAccent!5} پیش‌بینی: «حق شرعی» = ولایت \newline تطبیق: ● &
\cellcolor{cGold!8} پیش‌بینی: جهش فکری \newline تطبیق: ◐ &
\cellcolor{cConsolid!5} پیش‌بینی: حلقه تحمیل \newline تطبیق: ◐ \\
\midrule

\textbf{۳. قانون اساسی}{\footnotesize\\ (خرداد–آذر ۵۸)} &
تأیید اولیهٔ پیش‌نویس → چرخش به مجلس خبرگان → ولایت فقیه &
\cellcolor{cGreen!8} پیش‌بینی: حفظ پیش‌نویس \newline تطبیق: ○ &
\cellcolor{cAccent!5} پیش‌بینی: مهندسی ولایت فقیه \newline تطبیق: ● &
\cellcolor{cGold!8} پیش‌بینی: کشف ضرورت ولایت \newline تطبیق: ◐ &
\cellcolor{cConsolid!5} پیش‌بینی: بهشتی مهندسی کرد \newline تطبیق: ● \\
\midrule

\textbf{۴. فرقان}{\footnotesize\\ (فروردین–آذر ۵۸)} &
فرمان خلع سلاح، تقویت سپاه، توقیف مطبوعات &
\cellcolor{cGreen!8} پیش‌بینی: واکنش دفاعی \newline تطبیق: ◐ &
\cellcolor{cAccent!5} پیش‌بینی: بهانه برای سرکوب \newline تطبیق: ● &
\cellcolor{cGold!8} پیش‌بینی: تقویت باور به نیاز قدرت \newline تطبیق: ◐ &
\cellcolor{cConsolid!5} پیش‌بینی: سپاه بهره‌بردار \newline تطبیق: ● \\
\midrule

\textbf{۵. گروگان‌گیری}{\footnotesize\\ (آبان ۵۸)} &
تأیید ظرف ۲۴ ساعت، حذف بازرگان، تسریع رفراندوم &
\cellcolor{cGreen!8} پیش‌بینی: ترس از کودتا \newline تطبیق: ◐ &
\cellcolor{cAccent!5} پیش‌بینی: هماهنگ‌شده \newline تطبیق: ● &
\cellcolor{cGold!8} پیش‌بینی: کشف قدرت بسیج \newline تطبیق: ◐ &
\cellcolor{cConsolid!5} پیش‌بینی: ابتکار حزب ج.ا. \newline تطبیق: ● \\
\midrule

\textbf{۶. جنگ}{\footnotesize\\ (شهریور ۵۹–۶۷)} &
فرمانده‌کل، ادامهٔ جنگ پس از خرمشهر، هر نقد = خیانت &
\cellcolor{cGreen!8} پیش‌بینی: دفاع اضطراری \newline تطبیق: ◐ (اول) ○ (ادامه) &
\cellcolor{cAccent!5} پیش‌بینی: ابزار تثبیت \newline تطبیق: ● &
\cellcolor{cGold!8} پیش‌بینی: «مأموریت الهی» \newline تطبیق: ◐ &
\cellcolor{cConsolid!5} پیش‌بینی: سپاه ذی‌نفع \newline تطبیق: ● \\
\midrule

\textbf{۷. تصفیه‌ها}{\footnotesize\\ (۱۳۶۰–۶۸)} &
عزل بنی‌صدر، اعدام‌ها، کشتار ۶۷، عزل منتظری &
\cellcolor{cGreen!8} پیش‌بینی: واکنش به ترور \newline تطبیق: ○ &
\cellcolor{cAccent!5} پیش‌بینی: حذف سیستماتیک \newline تطبیق: ● &
\cellcolor{cGold!8} پیش‌بینی: «رحم = ضعف» \newline تطبیق: ◐ &
\cellcolor{cConsolid!5} پیش‌بینی: رقابت جناحی \newline تطبیق: ● \\

\bottomrule
\caption{ماتریس تطبیقی اصلی: هفت بزنگاه × چهار تز (● بالا ◐ متوسط ○ پایین)}
\label{tab:master-matrix}
\end{longtable}

\end{landscape}

% ────────────────────────────────────────
\section{شمارش امتیازها}
% ────────────────────────────────────────

\begin{principleBox}[جدول امتیازات]
اگر ● = ۲، ◐ = ۱، ○ = ۰ در نظر بگیریم:

\begin{longtable}{%
    >{\centering\arraybackslash}p{0.15\textwidth}
    >{\centering\arraybackslash}p{0.18\textwidth}
    >{\centering\arraybackslash}p{0.18\textwidth}
    >{\centering\arraybackslash}p{0.18\textwidth}
    >{\centering\arraybackslash}p{0.18\textwidth}}
\toprule
\rowcolor{cPrimary!10}
& \cellcolor{cGreen!15}\textbf{تز ۱} & \cellcolor{cAccent!10}\textbf{تز ۲} & \cellcolor{cGold!15}\textbf{تز ۳} & \cellcolor{cConsolid!10}\textbf{تز ۴} \\
\midrule

\textbf{مجموع ●} & ۱ & ۷ & ۱ & ۴ \\
\textbf{مجموع ◐} & ۴ & ۰ & ۶ & ۲ \\
\textbf{مجموع ○} & ۲ & ۰ & ۰ & ۱ \\
\midrule
\textbf{امتیاز کل (از ۱۴)} &
\cellcolor{cGreen!15}\textbf{۶} &
\cellcolor{cAccent!10}\textbf{۱۴} &
\cellcolor{cGold!15}\textbf{۸} &
\cellcolor{cConsolid!10}\textbf{۱۰} \\

\bottomrule
\end{longtable}
\end{principleBox}

% ────────────────────────────────────────
\section{تحلیل الگوها}
% ────────────────────────────────────────

\subsection{الگوی ۱: جهت ثابت تمرکز قدرت}

\begin{warningBox}[یافتهٔ کلیدی]
در \keyword{تمام هفت} بزنگاه، بدون استثنا، خمینی گزینه‌ای را انتخاب کرد که به \keyword{تمرکز بیشتر قدرت} منجر شد:

\begin{itemize}[label=\textcolor{cAccent}{$\blacktriangleright$}, rightmargin=2em]
    \item وقتی گزینهٔ «بازگشت به قم» وجود داشت، نرفت.
    \item وقتی پیش‌نویس دموکراتیک آماده بود، کنارش گذاشت.
    \item وقتی فرصت صلح بود (خرمشهر)، جنگ را ادامه داد.
    \item وقتی رئیس‌جمهور منتخب مخالفت کرد، عزلش کرد.
    \item وقتی جانشینش اعتراض کرد، عزلش کرد.
\end{itemize}

\textbf{احتمال تصادفی}: اگر هر بزنگاه دو گزینه داشته باشد (تمرکز / عدم تمرکز)، احتمال اینکه ۷ بار متوالی «اتفاقاً» تمرکز انتخاب شود: $(0.5)^7 = 0.78\%$ — کمتر از ۱\%. این الگو \keyword{تصادفی نیست}.
\end{warningBox}

\subsection{الگوی ۲: بهانه‌های متغیر، نتیجهٔ ثابت}

\begin{noteBox}[تنوع بهانه‌ها]
هر بار بهانهٔ متفاوتی وجود داشت:
\begin{itemize}[label=\textcolor{cGold}{$\bullet$}, rightmargin=2em]
    \item ۱۳۵۷: «مردم خواستند» (مشروعیت مردمی)
    \item ۱۳۵۸ (قانون اساسی): «اسلام نیاز دارد» (مشروعیت شرعی)
    \item ۱۳۵۸ (فرقان): «امنیت تهدید شده» (اضطرار امنیتی)
    \item ۱۳۵۸ (گروگان): «آمریکا توطئه می‌کند» (تهدید خارجی)
    \item ۱۳۵۹ (جنگ): «دفاع مقدس» (جنگ)
    \item ۱۳۶۰ (تصفیه): «منافقین خائنند» (تروریسم)
    \item ۱۳۶۸ (منتظری): «ساده‌لوحی» (ناشایستگی)
\end{itemize}

بهانه‌ها \keyword{هفت‌بار عوض شدند} اما نتیجه \keyword{یک‌بار هم} عوض نشد.
\end{noteBox}

\subsection{الگوی ۳: سرعت فزایندهٔ تمرکز}

\begin{figure}[H]
\centering
\begin{tikzpicture}[
    x=1.6cm,
    y=0.8cm
]
% محورها
\draw[-{Stealth[length=4pt]}, cGray, line width=0.5pt] (0,0) -- (9,0) node[below, font=\footnotesize] {زمان};
\draw[-{Stealth[length=4pt]}, cGray, line width=0.5pt] (0,0) -- (0,9) node[right, font=\footnotesize] {میزان تمرکز قدرت};

% نقاط
\fill[cPrimary] (0.5,1) circle (3pt) node[above right, font=\scriptsize, text width=1.8cm, align=center] {پاریس\\«مقامی ندارم»};
\fill[cPrimary] (1.5,3) circle (3pt) node[above right, font=\scriptsize, text width=1.5cm, align=center] {بازگشت\\«دولت تعیین می‌کنم»};
\fill[cPrimary] (3,5) circle (3pt) node[above right, font=\scriptsize, text width=1.5cm, align=center] {ولایت فقیه\\در قانون اساسی};
\fill[cPrimary] (4,5.5) circle (3pt) node[below right, font=\scriptsize, text width=1.5cm, align=center] {گروگان‌گیری\\حذف لیبرال‌ها};
\fill[cPrimary] (5,6.5) circle (3pt) node[above right, font=\scriptsize, text width=1.5cm, align=center] {جنگ\\فرماندهی کل};
\fill[cPrimary] (6.5,7.5) circle (3pt) node[above right, font=\scriptsize, text width=1.5cm, align=center] {عزل بنی‌صدر\\اعدام‌ها};
\fill[cPrimary] (8,8.5) circle (3pt) node[above right, font=\scriptsize, text width=1.5cm, align=center] {ولایت مطلقه\\عزل منتظری};

% خط روند
\draw[cAccent, line width=2pt, smooth] plot coordinates {(0.5,1) (1.5,3) (3,5) (4,5.5) (5,6.5) (6.5,7.5) (8,8.5)};

% برچسب سال‌ها
\node[font=\scriptsize\color{cGray}] at (0.5,-0.5) {۵۷};
\node[font=\scriptsize\color{cGray}] at (1.5,-0.5) {بهمن ۵۷};
\node[font=\scriptsize\color{cGray}] at (3,-0.5) {آذر ۵۸};
\node[font=\scriptsize\color{cGray}] at (4,-0.5) {آبان ۵۸};
\node[font=\scriptsize\color{cGray}] at (5,-0.5) {۵۹};
\node[font=\scriptsize\color{cGray}] at (6.5,-0.5) {۶۰};
\node[font=\scriptsize\color{cGray}] at (8,-0.5) {۶۶–۶۸};

\end{tikzpicture}
\caption{منحنی تمرکز قدرت: از «مقامی ندارم» تا «ولایت مطلقه» (۱۳۵۷–۱۳۶۸)}
\label{fig:power-curve}
\end{figure}

\subsection{الگوی ۴: عدم تقارن در داوری}

\begin{warningBox}[داوری یک‌طرفه]
هر بار که خمینی در موقعیت «داور» قرار گرفت، داوری‌اش به سود \keyword{جناح تندروتر} بود:

\begin{longtable}{%
    >{\raggedleft\arraybackslash}p{0.25\textwidth}
    >{\raggedleft\arraybackslash}p{0.25\textwidth}
    >{\raggedleft\arraybackslash}p{0.15\textwidth}
    >{\raggedleft\arraybackslash}p{0.25\textwidth}}
\toprule
\rowcolor{cAccent!10}
\textbf{تعارض} & \textbf{طرف میانه‌رو} & \textbf{طرف تندرو} & \textbf{داوری خمینی} \\
\midrule

دولت موقت / حزب ج.ا. & بازرگان & بهشتی & ✗ بازرگان \\
\midrule
رئیس‌جمهور / مجلس & بنی‌صدر & رجایی & ✗ بنی‌صدر \\
\midrule
پیش‌نویس / مجلس خبرگان & حبیبی-کاتوزیان & بهشتی & ✗ پیش‌نویس \\
\midrule
صلح / ادامهٔ جنگ & بازرگان، بنی‌صدر & سپاه & ✗ صلح \\
\midrule
رأفت / سرکوب & منتظری & دادگاه‌ها & ✗ منتظری \\
\midrule
\multicolumn{3}{r}{\textbf{دفعات داوری به سود تندروها:}} & \textbf{۵ از ۵ = ۱۰۰\%} \\

\bottomrule
\end{longtable}
\end{warningBox}

% ────────────────────────────────────────
\section{نمودار راداری: عملکرد تزها در هر بزنگاه}
% ────────────────────────────────────────

\begin{figure}[H]
\centering
\begin{tikzpicture}[scale=1.1]
    % تعریف محورها (۷ بزنگاه)
    \def\n{7}
    \def\labels{{"پاریس","بازگشت","قانون اساسی","فرقان","گروگان‌گیری","جنگ","تصفیه‌ها"}}

    % شبکه
    \foreach \r in {0.5,1,1.5,2} {
        \draw[cGray!20, thin] (0,0)
            \foreach \i in {1,...,\n} {
                -- ({90-360/\n*(\i-1)}:\r)
            } -- cycle;
    }

    % محورها
    \foreach \i in {1,...,\n} {
        \draw[cGray!30, thin] (0,0) -- ({90-360/\n*(\i-1)}:2.3);
        \pgfmathsetmacro{\angle}{90-360/\n*(\i-1)}
        \node[font=\scriptsize, text width=1.5cm, align=center] at (\angle:2.8) {\pgfmathparse{\labels[\i-1]}\pgfmathresult};
    }

    % تز ۱ (سبز)
    \draw[cGreen, thick, fill=cGreen!15, fill opacity=0.3]
        ({90-0*360/7}:1) --
        ({90-1*360/7}:0.5) --
        ({90-2*360/7}:0) --
        ({90-3*360/7}:1) --
        ({90-4*360/7}:0.5) --
        ({90-5*360/7}:0.5) --
        ({90-6*360/7}:0) -- cycle;

    % تز ۲ (قرمز)
    \draw[cAccent, thick, fill=cAccent!15, fill opacity=0.3]
        ({90-0*360/7}:2) --
        ({90-1*360/7}:2) --
        ({90-2*360/7}:2) --
        ({90-3*360/7}:2) --
        ({90-4*360/7}:2) --
        ({90-5*360/7}:2) --
        ({90-6*360/7}:2) -- cycle;

    % تز ۳ (طلایی)
    \draw[cGold, thick, fill=cGold!15, fill opacity=0.3]
        ({90-0*360/7}:2) --
        ({90-1*360/7}:1) --
        ({90-2*360/7}:1) --
        ({90-3*360/7}:1) --
        ({90-4*360/7}:1) --
        ({90-5*360/7}:1) --
        ({90-6*360/7}:1) -- cycle;

    % تز ۴ (بنفش)
    \draw[cConsolid, thick, fill=cConsolid!15, fill opacity=0.3]
        ({90-0*360/7}:1) --
        ({90-1*360/7}:1) --
        ({90-2*360/7}:2) --
        ({90-3*360/7}:2) --
        ({90-4*360/7}:2) --
        ({90-5*360/7}:2) --
        ({90-6*360/7}:2) -- cycle;

    % راهنما
    \node[font=\scriptsize, cGreen] at (-3.5,2) {— تز ۱ (صداقت)};
    \node[font=\scriptsize, cAccent] at (-3.5,1.5) {— تز ۲ (خدعه)};
    \node[font=\scriptsize, cGold] at (-3.5,1) {— تز ۳ (تکامل)};
    \node[font=\scriptsize, cConsolid] at (-3.5,0.5) {— تز ۴ (نخبگان)};

\end{tikzpicture}
\caption{نمودار راداری: عملکرد هر تز در هفت بزنگاه (مساحت بیشتر = تطبیق بهتر)}
\label{fig:radar-theses}
\end{figure}

% ────────────────────────────────────────
\section{نتیجه‌گیری تطبیقی}
% ────────────────────────────────────────

\begin{principleBox}[یافته‌های فصل ۱۳]
\begin{enumerate}[label=\textcolor{cPrimary}{\bfseries\arabic*.}, rightmargin=2em]
    \item \textbf{تز ۲ (خدعه)} بالاترین تطبیق را با واقعیت دارد: ۱۴ از ۱۴ (= ۱۰۰\%). در \keyword{هر} بزنگاه، رفتار خمینی دقیقاً همان بود که تز خدعه پیش‌بینی می‌کرد.
    \item \textbf{تز ۴ (نخبگان)} دومین رتبه: ۱۰ از ۱۴ (≈ ۷۱\%). در بیشتر بزنگاه‌ها (به‌ویژه قانون اساسی، فرقان، گروگان‌گیری، جنگ، تصفیه‌ها) نقش حلقهٔ درونی آشکار بود.
    \item \textbf{تز ۳ (تکامل)} سومین رتبه: ۸ از ۱۴ (≈ ۵۷\%). تبیین تکاملی در برخی بزنگاه‌ها (پاریس، ترور مطهری) قانع‌کننده بود اما در بزنگاه‌های حاد (تصفیه‌ها، کشتار ۶۷) ضعیف عمل کرد.
    \item \textbf{تز ۱ (صداقت)} ضعیف‌ترین: ۶ از ۱۴ (≈ ۴۳\%). فقط در بزنگاه پاریس (خود وعده‌ها) تطبیق بالایی داشت. در بقیهٔ بزنگاه‌ها ناتوان از توضیح جهت ثابت تمرکز بود.
    \item \keyword{الگوی کلی:} جهت ثابت تمرکز، بهانه‌های متغیر، سرعت فزاینده، داوری یکسویه.
\end{enumerate}
\end{principleBox}


% ============================================================
%  فصل ۱۴: خوانش‌های رقیب در بوتهٔ نقد
% ============================================================
\chapter{خوانش‌های رقیب در بوتهٔ نقد}

\begin{infoBox}[چکیدهٔ فصل]
پس از تحلیل تطبیقی فصل ۱۳، اکنون هر تز را به‌طور سیستماتیک از منظر سه تز دیگر نقد می‌کنیم. شش جفت (\en{pair}) نقد متقابل بررسی می‌شود. سپس پرسیده می‌شود: آیا تزها \keyword{ترکیب‌پذیر} هستند یا متعارض؟ ماتریس سازگاری ارائه می‌شود.
\end{infoBox}

% ────────────────────────────────────────
\section{شش جفت نقد متقابل}
% ────────────────────────────────────────

\subsection{جفت ۱: تز ۱ (صداقت) ↔ تز ۲ (خدعه)}

\begin{longtable}{%
    >{\raggedleft\arraybackslash}p{0.45\textwidth}
    >{\raggedleft\arraybackslash}p{0.45\textwidth}}
\toprule
\rowcolor{cGreen!10}
\textbf{تز ۱ علیه تز ۲} &
\rowcolor{cAccent!8}
\textbf{تز ۲ علیه تز ۱} \\
\midrule

فریبکار هوشمند ابهام می‌گذارد نه وعدهٔ صریح. خمینی ۱۷+ بار صریحاً وعده داد — این با فریب ناسازگار است. &
کسی که ده سال قبل کتاب «حکومت اسلامی» نوشته نمی‌تواند صادقانه بگوید «حکومت نخواهم کرد» — مگر تعریفش از «حکومت» متفاوت باشد (= بازی زبانی = فریب). \\
\midrule

بازرگان (غیرروحانی) نخست‌وزیر شد — فریبکار از ابتدا روحانی می‌گذاشت. &
بازرگان «پوشش» بود. خمینی هم‌زمان شورای انقلاب مخفی داشت — دولت موقت صوری بود. \\
\midrule

فرض «طراحی بلندمدت» نیازمند ذهن تمام‌بین است. واقعیت آشفته‌تر از توطئه‌هاست. &
نیاز به «ذهن تمام‌بین» نیست — فقط به «جهت ثابت» نیاز هست. خمینی نباید تمام مسیر را از قبل بداند؛ کافی‌ست در هر بزنگاه «تمرکز بیشتر» را انتخاب کند. \\

\bottomrule
\caption{نقد متقابل: تز ۱ ↔ تز ۲}
\end{longtable}

\begin{noteBox}[ارزیابی جفت ۱]
نقد تز ۲ علیه تز ۱ \keyword{قوی‌تر} است: وجود کتاب «حکومت اسلامی» (۱۳۴۸) سنگین‌ترین شاهد علیه صداقت محض است. پاسخ تز ۱ (ثبات وعده‌ها = صداقت) تا حدی معتبر است اما نمی‌تواند فاصلهٔ ساختاری وعده و عمل را توضیح دهد.
\end{noteBox}

\subsection{جفت ۲: تز ۱ (صداقت) ↔ تز ۳ (تکامل)}

\begin{longtable}{%
    >{\raggedleft\arraybackslash}p{0.45\textwidth}
    >{\raggedleft\arraybackslash}p{0.45\textwidth}}
\toprule
\rowcolor{cGreen!10}
\textbf{تز ۱ علیه تز ۳} &
\rowcolor{cGold!10}
\textbf{تز ۳ علیه تز ۱} \\
\midrule

«تکامل فکری» فرض می‌کند خمینی تردید داشت — هیچ شاهد مستقیمی بر تردید وجود ندارد. &
تز ۱ خمینی را «واکنش‌گر منفعل» نشان می‌دهد — اما خمینی در هر بزنگاه فعالانه انتخاب کرد. «اجبار» مفهوم ضعیفی‌ست وقتی فرد گزینه دارد و همواره یک جهت را انتخاب می‌کند. \\
\midrule

سرعت تغییرات (۳۰ روز از پاریس تا بهشت‌زهرا) با «تکامل تدریجی» ناسازگار است — تکامل آهسته‌تر از این‌هاست. &
تکامل لزوماً خطی نیست — لحظات «جهش» وجود دارد. بازگشت به ایران و مواجهه با واقعیت قدرت ممکن بود «لحظهٔ جهش» باشد. \\

\bottomrule
\caption{نقد متقابل: تز ۱ ↔ تز ۳}
\end{longtable}

\subsection{جفت ۳: تز ۱ (صداقت) ↔ تز ۴ (نخبگان)}

\begin{longtable}{%
    >{\raggedleft\arraybackslash}p{0.45\textwidth}
    >{\raggedleft\arraybackslash}p{0.45\textwidth}}
\toprule
\rowcolor{cGreen!10}
\textbf{تز ۱ علیه تز ۴} &
\rowcolor{cConsolid!8}
\textbf{تز ۴ علیه تز ۱} \\
\midrule

تز ۴ مسئلهٔ «مرغ و تخم‌مرغ» دارد: نخبگان بدون خمینی قدرتی نداشتند. &
حتی اگر بپذیریم خمینی صادق بود، تأثیر نخبگان بر او غیرقابل‌انکار است. «صداقت» خمینی مانع تأثیرپذیری‌اش از حلقهٔ درونی نمی‌شود. \\
\midrule

تز ۴ عاملیت فردی خمینی را بیش از حد کم‌اهمیت می‌کند. &
تز ۱ نقش ساختارها و نهادها را بیش از حد نادیده می‌گیرد و همه‌چیز را به «اراده» یا «اجبار» یک فرد تقلیل می‌دهد. \\

\bottomrule
\caption{نقد متقابل: تز ۱ ↔ تز ۴}
\end{longtable}

\subsection{جفت ۴: تز ۲ (خدعه) ↔ تز ۳ (تکامل)}

\begin{longtable}{%
    >{\raggedleft\arraybackslash}p{0.45\textwidth}
    >{\raggedleft\arraybackslash}p{0.45\textwidth}}
\toprule
\rowcolor{cAccent!8}
\textbf{تز ۲ علیه تز ۳} &
\rowcolor{cGold!10}
\textbf{تز ۳ علیه تز ۲} \\
\midrule

«تکامل» توضیح ملایم‌تری است برای چیزی که ساده‌تر با «پنهان‌کاری» توضیح‌پذیر است (\en{Occam's Razor}). &
تز ۲ فرض «ذهن ثابت از ۱۳۴۸» دارد — اما متون خمینی نشان‌دهندهٔ تغییرات واقعی در مواضع است (کشف‌الاسرار ≠ حکومت اسلامی ≠ ولایت مطلقه). \\
\midrule

هیچ شاهد مستقیمی بر «تردید» خمینی نیست — فقط حدس. &
هیچ شاهد مستقیمی بر «طراحی کامل از ابتدا» هم نیست — فقط استنتاج از الگو. هر دو تز حدس‌محورند. \\

\bottomrule
\caption{نقد متقابل: تز ۲ ↔ تز ۳}
\end{longtable}

\subsection{جفت ۵: تز ۲ (خدعه) ↔ تز ۴ (نخبگان)}

\begin{longtable}{%
    >{\raggedleft\arraybackslash}p{0.45\textwidth}
    >{\raggedleft\arraybackslash}p{0.45\textwidth}}
\toprule
\rowcolor{cAccent!8}
\textbf{تز ۲ علیه تز ۴} &
\rowcolor{cConsolid!8}
\textbf{تز ۴ علیه تز ۲} \\
\midrule

تز ۴ نقش خمینی را کم‌اهمیت می‌کند. خمینی «بازیکن منفعل» نبود — او وتوی نهایی داشت و هیچ‌گاه استفاده نکرد. &
تز ۲ همه‌چیز را به یک فرد نسبت می‌دهد. واقعیت پیچیده‌تر بود: بهشتی، رفسنجانی، و سپاه هر‌یک «عاملیت مستقل» داشتند. \\
\midrule

اگر نخبگان «ابتکار عمل» داشتند، چرا همواره نتیجه به سود \textit{خمینی} تمام شد؟ &
چون خمینی «نقطهٔ ثقل» بود — هم نخبگان و هم خمینی سود بردند. رابطه \keyword{همافزایی} بود نه تبعیت. \\

\bottomrule
\caption{نقد متقابل: تز ۲ ↔ تز ۴}
\end{longtable}

\subsection{جفت ۶: تز ۳ (تکامل) ↔ تز ۴ (نخبگان)}

\begin{longtable}{%
    >{\raggedleft\arraybackslash}p{0.45\textwidth}
    >{\raggedleft\arraybackslash}p{0.45\textwidth}}
\toprule
\rowcolor{cGold!10}
\textbf{تز ۳ علیه تز ۴} &
\rowcolor{cConsolid!8}
\textbf{تز ۴ علیه تز ۳} \\
\midrule

تز ۴ بُعد فکری را نادیده می‌گیرد. ولایت فقیه صرفاً «محصول رقابت نخبگان» نبود — مبنای فقهی داشت. &
مبنای فقهی ابزار بود، نه علت. اگر ولایت فقیه به نفع نخبگان نبود، هیچ‌کس آن‌را «کشف» نمی‌کرد. \\
\midrule

تکامل فکری مستقل از فشار نخبگان بود — خمینی از قبل نظریه داشت. &
داشتن نظریه یک چیز است و اجرای آن چیز دیگر. نخبگان بستر اجرا را فراهم کردند. \\

\bottomrule
\caption{نقد متقابل: تز ۳ ↔ تز ۴}
\end{longtable}

% ────────────────────────────────────────
\section{ماتریس سازگاری: آیا تزها ترکیب‌پذیرند؟}
% ────────────────────────────────────────

\begin{principleBox}[ماتریس سازگاری]
\begin{longtable}{%
    >{\centering\arraybackslash}p{0.15\textwidth}
    >{\centering\arraybackslash}p{0.18\textwidth}
    >{\centering\arraybackslash}p{0.18\textwidth}
    >{\centering\arraybackslash}p{0.18\textwidth}
    >{\centering\arraybackslash}p{0.18\textwidth}}
\toprule
\rowcolor{cPrimary!10}
& \cellcolor{cGreen!15}\textbf{تز ۱} & \cellcolor{cAccent!10}\textbf{تز ۲} & \cellcolor{cGold!15}\textbf{تز ۳} & \cellcolor{cConsolid!10}\textbf{تز ۴} \\
\midrule

\cellcolor{cGreen!15}\textbf{تز ۱} & — & \textcolor{cAccent}{ناسازگار} & \textcolor{cGold}{نیمه‌سازگار} & \textcolor{cGreen}{سازگار} \\
\midrule
\cellcolor{cAccent!10}\textbf{تز ۲} & \textcolor{cAccent}{ناسازگار} & — & \textcolor{cGold}{نیمه‌سازگار} & \textcolor{cGreen}{سازگار} \\
\midrule
\cellcolor{cGold!15}\textbf{تز ۳} & \textcolor{cGold}{نیمه‌سازگار} & \textcolor{cGold}{نیمه‌سازگار} & — & \textcolor{cGreen}{سازگار} \\
\midrule
\cellcolor{cConsolid!10}\textbf{تز ۴} & \textcolor{cGreen}{سازگار} & \textcolor{cGreen}{سازگار} & \textcolor{cGreen}{سازگار} & — \\

\bottomrule
\end{longtable}

\textbf{خوانش ماتریس:}
\begin{itemize}[label=\textcolor{cPrimary}{$\bullet$}, rightmargin=2em]
    \item \textcolor{cGreen}{\textbf{سازگار:}} دو تز می‌توانند هم‌زمان درست باشند.
    \item \textcolor{cGold}{\textbf{نیمه‌سازگار:}} در بعضی جنبه‌ها قابل ترکیب، در بعضی متعارض.
    \item \textcolor{cAccent}{\textbf{ناسازگار:}} متعارض ذاتی — اگر یکی درست باشد، دیگری غلط است.
\end{itemize}
\end{principleBox}

\begin{noteBox}[مهم‌ترین یافته]
\keyword{تز ۴ (نخبگان) با هر سه تز دیگر سازگار است} — یعنی «فشار نخبگان» یک \keyword{مکمل عمومی} است نه رقیب. این تز فراگیرترین چارچوب را دارد.

\keyword{تز ۱ و تز ۲ ناسازگار ذاتی} هستند: یا صادق بود یا فریب داد — هر دو هم‌زمان نمی‌شود (مگر بپذیریم بخشی صادقانه و بخشی فریبکارانه بود — که ما را به تز ۳ می‌رساند).

\keyword{بهترین ترکیب:} تز ۲ + تز ۳ + تز ۴ = خمینی احتمالاً \textit{هم} استراتژیک بود، \textit{هم} تکامل فکری داشت، و \textit{هم} تحت تأثیر نخبگان بود.
\end{noteBox}

% ────────────────────────────────────────
\section{جدول نقاط قوت و ضعف هر تز}
% ────────────────────────────────────────

\begin{longtable}{%
    >{\centering\arraybackslash}p{0.07\textwidth}
    >{\raggedleft\arraybackslash}p{0.14\textwidth}
    >{\raggedleft\arraybackslash}p{0.33\textwidth}
    >{\raggedleft\arraybackslash}p{0.33\textwidth}}
\toprule
\rowcolor{cPrimary!10}
\textbf{تز} & \textbf{نام} & \textbf{نقاط قوت} & \textbf{نقاط ضعف} \\
\midrule
\endhead

\cellcolor{cGreen!10} ۱ &
\cellcolor{cGreen!10} صداقت + اجبار &
\cellcolor{cGreen!10} ثبات وعده‌ها در پاریس؛ انتصاب بازرگان؛ بحران‌های واقعی وجود داشت &
\cellcolor{cGreen!10} \keyword{ناتوان} از توضیح جهت ثابت تمرکز؛ نادیده‌گرفتن کتاب حکومت اسلامی؛ تقلیل خمینی به واکنش‌گر \\
\midrule

\cellcolor{cAccent!5} ۲ &
\cellcolor{cAccent!5} خدعهٔ آگاهانه &
\cellcolor{cAccent!5} \keyword{قوی‌ترین} تطبیق با واقعیت؛ توضیح فاصلهٔ عمیق وعده/عمل؛ وجود «حکومت اسلامی» ۱۳۴۸ &
\cellcolor{cAccent!5} فرض ذهن استراتژیک کامل؛ نادیده‌گرفتن نقش نخبگان؛ اتهام «توطئه‌اندیشی» \\
\midrule

\cellcolor{cGold!8} ۳ &
\cellcolor{cGold!8} تکامل تدریجی &
\cellcolor{cGold!8} توضیح تغییرات مستمر در متون خمینی؛ انعطاف‌پذیری؛ پرهیز از ساده‌سازی &
\cellcolor{cGold!8} ابهام ذاتی؛ فقدان شاهد مستقیم بر «تردید»؛ استقلال ضعیف (بیشتر مکمل) \\
\midrule

\cellcolor{cConsolid!5} ۴ &
\cellcolor{cConsolid!5} فشار نخبگان &
\cellcolor{cConsolid!5} تبیین ساختاری؛ سازگار با هر تز دیگر؛ توضیح نقش بازیگران متعدد &
\cellcolor{cConsolid!5} کم‌اهمیت‌نمایی عاملیت خمینی؛ مشکل مرغ‌وتخم‌مرغ؛ نادیده‌گرفتن بُعد فکری \\

\bottomrule
\caption{جدول تطبیقی نقاط قوت و ضعف چهار تز}
\label{tab:strengths-weaknesses}
\end{longtable}

% ────────────────────────────────────────
\section{جمع‌بندی فصل چهاردهم}
% ────────────────────────────────────────

\begin{principleBox}[خلاصهٔ فصل ۱۴]
\begin{itemize}[label=\textcolor{cPrimary}{$\blacksquare$}, rightmargin=2em]
    \item شش جفت نقد متقابل نشان داد هیچ تزی بدون نقص نیست.
    \item تز ۱ و تز ۲ ذاتاً ناسازگارند — باید یکی را (عمدتاً) انتخاب کرد.
    \item تز ۴ (نخبگان) با هر سه تز سازگار است و فراگیرترین چارچوب را دارد.
    \item تز ۳ (تکامل) نیمه‌سازگار با هر دو تز ۱ و ۲ است — می‌تواند «پل» باشد.
    \item \keyword{بهترین ترکیب پیشنهادی:} تز ۲ (خدعه/استراتژی) + تز ۳ (تکامل) + تز ۴ (نخبگان).
    \item فصل ۱۵ با ابزارهای روان‌شناختی این ترکیب را دقیق‌تر بررسی خواهد کرد.
\end{itemize}
\end{principleBox}